% Junior Level 3, Week 4 -- How Do Computers Learn? (unplugged: no computer needed this week)
\lesson{4}{How Do Computers Learn?}{Computers learn patterns from examples, and copy our mistakes.}
\lessoninfo{Concepts: training examples (I), features (I), predictions (I), bias (I), tally and bar graph (R)}{Tools: printed animal cards (page at the end of this plan), crayons; no computer needed}{CSTA: 1B-DA-07, 1B-IC-18, 1B-NI-05; AI4K12 Big Ideas 3 and 5}

\begin{objectives}
  \item The student can teach a partner to sort animal cards from examples and say which features, such as legs, wings or fur, helped. (J3.G4) [J3.W04.1]
  \item The student can test the \enquote{robot} with new cards, tally its right and wrong predictions, and show them in a bar graph. (J3.G3, J3.G4) [J3.W04.2]
  \item The student can explain how too few or one-sided examples lead to wrong guesses, and fix it by adding different examples. (J3.G4, J3.G6) [J3.W04.3]
  \item The student can say that AI can be wrong and that we never tell a chatbot personal information. (J3.G7) [J3.W04.4]
\end{objectives}

\begin{materials}
  \item Print: workbook chapter 4 and play-at-home card 4 (one per child), and the \textbf{animal card page} at the end of this plan, one copy per pair. Cut the cards apart (the children can help).
  \item Two sheets of paper labeled \textbf{YES} and \textbf{NO}, per pair.
  \item For the warm-up: 8 to 10 small things from around the house, some with wheels and some without (toy cars, a ball, a spoon, a toy train, a book\ldots).
  \item A kitchen timer for role swaps (6--7 minutes).
  \item No computer this week. Optional: the Rose and Pixi video for this week.
\end{materials}

\needspace{16\baselineskip}
\section{At a glance}
\begin{timeplan}
  \timeslot{0--10}{Warm-up}{Guess my rule: you sort things into YES and NO piles by a secret rule; the children guess the rule and predict where new things go.}
  \timeslot{10--20}{Show me}{How computers learn: examples, features and predictions. A demo round of Teach the Robot with you as the Trainer.}
  \timeslot{20--33}{Make it}{Exercise 4.1 Teach the Robot in pairs: Trainer and Robot, swap on the timer. Tally right and wrong guesses.}
  \timeslot{33--35}{Wiggle break}{Robot sort: four legs, stomp; can fly, flap; lives in water, swim.}
  \timeslot{35--50}{Make it}{Exercise 4.2 Robot Report (bar graph), Exercise 4.3 The One-Sided Robot (bug hunt), then free build: invent a creature card.}
  \timeslot{50--60}{Share}{What did the robot learn? AI can be wrong; never tell a chatbot personal information. Exit question, hand out play-at-home card 4.}
\end{timeplan}

\section{Warm-up (0--10 min)}
\begin{unplugged}[Guess my rule]
\begin{enumerate}[leftmargin=1.4em,itemsep=1pt]
  \item Secretly pick a rule: \enquote{has wheels}. Put the two paper labels YES and NO on the floor.
  \item Place things one at a time: a toy car on YES, a spoon on NO, a toy train on YES, a ball on NO. Say only \enquote{yes} or \enquote{no}, never the rule.
  \item Hold up a new thing: \enquote{Where does this go?} Children \textbf{predict}. Then show the answer. After 5 or 6 things, ask: \enquote{What is my rule?}
  \item \enquote{You learned my rule from \textbf{examples}, without me telling you. That is how many computers learn too.}
\end{enumerate}
With 1 child: play two rounds, then let the child pick a rule for you. With 3 or 4: children take turns picking the rule after the first round.
\end{unplugged}

\section{Show me (10--20 min)}
\begin{teacherscript}
\enquote{Some computer programs are not built block by block. Instead, people show them lots of \textbf{examples}, each one marked YES or NO. The computer looks for a pattern in the examples. This is called \textbf{machine learning}, and programs that learn like this are often called AI. The things the computer looks at, like legs, wings or fur, are called \textbf{features}. When it sees something new, it makes a guess, called a \textbf{prediction}. Today one of you is the computer: the Robot.}
\end{teacherscript}

\begin{analogy}[Learning what a dog is]
\enquote{Nobody gave you a rule list for \enquote{dog} when you were little. You saw lots of dogs, big and small, and someone said \enquote{dog!} each time. Then you could spot a dog you had never seen. A computer learns from examples in the same way, but it only knows the examples we give it.}
\end{analogy}

\begin{workedexample}[A demo round of Teach the Robot]
You are the Trainer; one child is the Robot. The others watch.
\begin{enumerate}[leftmargin=1.4em,itemsep=1pt]
  \item Secretly take the rule card \enquote{HAS FOUR LEGS}. Shuffle the animal cards.
  \item Show 6 \textbf{training examples}, one at a time, and put each on YES or NO: cat (YES), fish (NO), horse (YES), dove (NO), frog (YES), spider (NO).
  \item Now the \textbf{test}: hold up a new card, the dog. \enquote{Robot, YES or NO?} The Robot predicts. Say \enquote{right} or \enquote{wrong} and put the card on the correct pile. Do the same with the hippo and the crow.
  \item Ask the Robot: \enquote{What do you think the rule is? Which feature did you look at?} (Legs.) Show the rule card.
  \item \enquote{The Robot never saw the rule. It learned it from the examples, and then it could guess new ones.}
\end{enumerate}
\end{workedexample}

\section{Make it (20--33 min)}
\begin{pairroles}[Trainer and Robot]
This week the roles are \textbf{Trainer} and \textbf{Robot}. The Trainer picks a secret rule card, shows 6 examples on YES and NO, and then tests the Robot with new cards. The Robot sorts without seeing the rule, predicts, and at the end says the rule and the feature that helped. The Trainer only says \enquote{yes}, \enquote{no}, \enquote{right} or \enquote{wrong}: no hints. Swap when the timer rings (6--7 minutes). With 1 child: you are the Robot first (make a few honest mistakes), then swap. With 3: Trainer, Robot and Recorder (keeps the tally of right and wrong). With 4: two pairs; then pairs swap rule cards.
\end{pairroles}

Work through \textbf{4.1} Teach the Robot in the workbook. The Robot's partner keeps a tally of right and wrong predictions in the workbook. Rule cards: \enquote{HAS FOUR LEGS}, \enquote{IS A BIRD}, \enquote{CAN FLY}. Some cards are tricky on purpose: the kiwi is a bird that cannot fly, and the bug can fly but is not a bird.

\section{Wiggle break (33--35 min)}
Robot sort with the whole body. Call out an animal: if it has four legs, stomp; if it can fly, flap your arms; if it lives in water, make swimming arms. Some animals need two moves (a duck flaps and swims!).

\section{Make it (35--50 min)}
\textbf{4.2} Robot Report: color a bar graph of the Robot's right and wrong predictions from the tally. Ask: \enquote{Did the Robot get better after more examples?}

\begin{bughunt}[The One-Sided Robot (Exercise 4.3)]
Sam trained the Robot with the rule \enquote{IS A BIRD}, but every YES example was a bird that flies: dove, crow, dove. Every NO example was an animal without wings. Play it: one child is Sam's Robot and sees only those examples. Then test with the kiwi and the bug. The Robot says the kiwi is NOT a bird (it does not fly) and the bug IS a bird (it flies). Ask: \enquote{Is the Robot silly, or were the examples one-sided?} The Robot learned exactly what it was shown: \enquote{bird means flying}. Fix: add different examples: the kiwi on YES, the bug on NO. Train again and test again.

\enquote{Real computers make this mistake too. If an AI only learns from pictures of some kinds of people, or only hears some kinds of voices, it makes more mistakes with everyone else. That is unfair, and it is called \textbf{bias}. The fix is the same as ours: more and different examples, and people checking the results.}
\end{bughunt}

\begin{freebuild}[Invent a creature]
Last 6--7 minutes: on the blank card, draw a made-up creature with clear features (How many legs? Wings? Fur? Fins?). Give it a fun name. Then ask the Robot: \enquote{With the rule HAS FOUR LEGS, is my creature YES or NO?} Does the Robot agree with you?
\end{freebuild}

\begin{stretchchallenge}[A harder rule]
Pick a rule with two features: \enquote{has four legs AND has fur}. How many examples does the Robot need before it guesses right? Is the frog a YES or a NO now?
\end{stretchchallenge}

\section{Share (50--60 min)}
\begin{teacherscript}
Each Trainer says which rule they taught and holds up the Robot Report bar graph: \enquote{My Robot got 4 right and 1 wrong.} Each Robot says which feature helped. Others give a thumbs up and say one thing they like.
\end{teacherscript}

\begin{safety}[AI can be wrong, and it is not a friend]
\enquote{Some apps and websites have an AI chatbot: you type, and it writes back like a person. It learned from huge numbers of examples, so it can be very helpful, but, like our Robot, it can be wrong, even when it sounds very sure. If an answer matters, check it with a grown-up or a good book.}

\enquote{A chatbot is a computer program, not a friend. Never tell a chatbot personal information: not your name, your address, your school, your age, a photo of you, or a password. What you type might be kept. Only use a chatbot with a grown-up next to you.} Ask: \enquote{A chatbot asks: \enquote{What is your name and where do you live?} What do you do?} (Do not answer, close it, tell a grown-up.)
\end{safety}

\begin{formativecheck}[Exit question]
\textbf{Look for:} the Trainer says only yes/no and right/wrong; the Robot names a feature, not only a card; the tally and the bar graph show the same numbers.

\textbf{Ask:} \enquote{Omar taught a robot what a cat is, but he only showed it orange cats. The robot sees a black cat. What might it guess, and how can Omar fix it?}
\textbf{Answer:} it might guess \enquote{not a cat}, because every cat it saw was orange. Omar should show it cats of many colors.
\end{formativecheck}

\begin{reflection}
Ask: \enquote{How did the Robot learn?} (from examples) and \enquote{Why did Sam's Robot make mistakes?} (the examples were one-sided). Children draw their made-up creature in the workbook and circle a smiley face. Hand out play-at-home card 4.
\end{reflection}

\section{Notes for the teacher}
\begin{commonmistakes}
  \item \textbf{The Trainer gives hints} (\enquote{look at the legs!}). Then the Robot is told, not trained. Keep to yes, no, right, wrong.
  \item \textbf{The Robot only remembers cards} (\enquote{the cat was YES}). Ask: \enquote{What do the YES cards have in common?} Push toward features.
  \item \textbf{Arguing about a card} (Does a hippo live in water? Is a frog four-legged?). Great: real data is messy. Agree on one answer and say so; it is a good example of why people check AI.
  \item \textbf{\enquote{The robot is stupid.}} Turn it round: the robot learned what it was shown. Mistakes usually come from the examples.
  \item \textbf{Bar graph bars that do not match the tally.} Count the tally marks again, one square each.
\end{commonmistakes}

\begin{differentiation}
\supportpath Use only the rule \enquote{CAN FLY} and 4 training examples. You are the Trainer every round; the child is the Robot and you do the tally together.
\extensionpath The two-feature rule in the stretch challenge. Or: \enquote{What is the smallest number of examples that still teaches the rule?} Try 2, then 4, then 6.
\end{differentiation}

\begin{teachernote}[If you are short on time]
Skip the warm-up objects and go straight to the demo round. Do one round of 4.1 per child and the bug hunt (4.3); color the bar graph (4.2) at home. Never skip the chatbot safety talk.
\end{teachernote}

\begin{teachernote}[Background for the grown-up]
Today's activity is a small model of supervised machine learning: labeled training examples, features, predictions on new data, and checking accuracy. The bias part (one-sided examples lead to unfair mistakes) follows AI4K12 Big Idea 3 (computers can learn from data) and Big Idea 5 (AI can affect people, in good and bad ways). You do not need to use these words with the children.
\end{teachernote}

\clearpage
\section{Printable animal cards}
Print one copy per pair and cut along the dashed lines: 12 animal cards, 3 rule cards and a blank card for the free build. Every card has a picture and a word.

\begin{center}
%% alt: Sixteen cut-out cards with dashed borders. Twelve animal cards, each with a picture and a name: cat, dog, horse, hippo, otter, frog, fish, dove, crow, kiwi bird, spider and bug. Three rule cards with a lock: HAS FOUR LEGS, IS A BIRD, CAN FLY. One blank card that says draw your own creature.
\begin{tikzpicture}[x=1cm,y=1cm,
  edge/.style={draw=kcblack,dashed,line width=0.8pt,rounded corners=4pt},
  word/.style={font=\bfseries\large,align=center}]
  \foreach \icon/\name [count=\i from 0] in {%
      \faCat/cat, \faDog/dog, \faHorse/horse, \faHippo/hippo,
      \faOtter/otter, \faFrog/frog, \faFish/fish, \faDove/dove,
      \faCrow/crow, \faKiwiBird/kiwi bird, \faSpider/spider, \faBug/bug} {
    \pgfmathtruncatemacro{\col}{mod(\i,4)}
    \pgfmathtruncatemacro{\row}{3-div(\i,4)}
    \begin{scope}[shift={(\col*4.2,\row*4.2)}]
      \filldraw[edge,fill=white] (0,0) rectangle (3.8,3.8);
      \node[scale=3] at (1.9,2.3) {\icon};
      \node[word] at (1.9,0.6) {\name};
    \end{scope}}
  \foreach \rulename [count=\i from 0] in {HAS FOUR\\LEGS, IS A\\BIRD, CAN\\FLY} {
    \begin{scope}[shift={(\i*4.2,0)}]
      \filldraw[edge,fill=kcyellow!25] (0,0) rectangle (3.8,3.8);
      \node[scale=1.6] at (1.9,2.9) {\faLock};
      \node[word] at (1.9,1.4) {\rulename};
    \end{scope}}
  \begin{scope}[shift={(12.6,0)}]
    \filldraw[edge,fill=kcblue!8] (0,0) rectangle (3.8,3.8);
    \node[word,font=\bfseries] at (1.9,0.7) {draw your\\own creature};
  \end{scope}
\end{tikzpicture}
\end{center}
